% ---------------------------------------------------------------------------
\section{Single-Problem Self-Reinforcement}
\label{sec:single}

We start with one problem in isolation. Its sharpening drift $F$, and in particular the slope $F'$, are the only single-problem ingredients that the collective theory of \cref{sec:theory} needs.

\paragraph{Policy, reward and drift.}
We abstract a problem $x$ by the distribution of its \emph{final answers}.
In the binary case the two answers are $+$ and $-$, and the policy is $\pi(+\mid x)=p=\sigma(z)$, $\pi(-\mid x)=1-p$, with logit $z$.
For each problem the learner samples $K$ answers $y_1,\dots,y_K$, forms the majority $\hat a$ (ties broken uniformly at random), and assigns $r_j=\mathbf 1[y_j=\hat a]$ \citep{ttrl}.
The per-sample policy-gradient estimate is $\hat g=\frac1K\sum_j A_j\,\partial_z\log\pi(y_j)$, with $\partial_z\log\pi(+)=1-p$, $\partial_z\log\pi(-)=-p$, and advantage $A_j=r_j$ (REINFORCE; \citealp{williams1992}), $A_j=r_j-\bar r$ (group-mean baseline), or $A_j=(r_j-\bar r)/\mathrm{std}(r)$ (GRPO; \citealp{deepseekmath}).

The estimate $\hat g$ averages over the $K$ answers drawn for a single update, so it is still random: each update draws fresh answers. Its expectation over the draw of $y_1,\dots,y_K\sim\pi$ and the tie-breaking is the mean update on the logit, $F(z)=\mathbb E[\hat g]$.
Because $\hat g$ depends on the samples only through the number $k$ of $+$ answers, $F$ can be computed exactly by summing over $k\sim\mathrm{Bin}(K,p)$ for any $K$ and estimator (\cref{app:drift}).
For REINFORCE it has a simple form. By exchangeability $F=\mathbb E[r_1\,\partial_z\log\pi(y_1)]$, so $F(z)=p(1-p)\big(\Pr[\hat a=+\mid y_1=+]-\Pr[\hat a=-\mid y_1=-]\big)$.
For odd $K=2m+1$ the bracket equals $\Pr[X>m]-\Pr[X<m]$, where $X\sim\mathrm{Bin}(2m,p)$ counts the $+$ answers among the other $K-1$ samples.
Our main experiments use $K=3$. There the bracket is $p^2-(1-p)^2=2p-1$, which gives
\begin{equation}
  F(z)=p(1-p)(2p-1),\qquad p=\sigma(z).
  \label{eq:F}
\end{equation}
A problem trained on its own with step $\eta$ therefore follows, on average, the \emph{sharpening equation}
\begin{equation}
  \dot z=\eta\,F(z),
  \label{eq:sharp}
\end{equation}
where time counts updates.
For every $K\ge3$ and all three estimators, $F$ is odd, positive for $z>0$, and vanishes at $z=0$ and as $|z|\to\infty$. Under \cref{eq:sharp} every problem is therefore pushed toward its current majority answer and then freezes.
This is the single-problem sharpening behaviour characterized by \citet{huang2025sharpening} and \citet{he2026urlvr}.

\paragraph{The unstable window.}
The drift $F$ describes how a problem sharpens on its own. Once problems share parameters, they also feel a common push, and how a problem responds to that push is set by the \emph{slope} $F'(z)$. This slope is the quantity that controls the collective behaviour.
For $K=3$, differentiating \cref{eq:F} and writing $q=p-\tfrac12$ gives $F'(z)=p(1-p)\,(\tfrac12-6q^2)$. Setting $6q^2=\tfrac12$ locates the zero of $F'$, so
\begin{equation}
  F'(z)>0 \iff |z|<z^\star=\ln(2+\sqrt3)\approx1.317,
  \label{eq:window}
\end{equation}
that is, $0.21<p<0.79$. Evaluating \cref{eq:F} at the edge gives $F^\star=F(z^\star)=1/(6\sqrt3)\approx0.096$.
Inside this window a problem's response to a small common push grows with the push. Outside the window it shrinks (\cref{fig:mechanism}a).

So far we have used REINFORCE with $K=3$. We now check how $F$ changes with the estimator and the group size.
For group-mean and GRPO advantages, unanimous groups have zero advantage, so only the $2{:}1$ splits contribute, and they all have the same reward standard deviation. For REINFORCE, unanimous groups also contribute; in fact for $K=3$ they carry the entire mean drift, while the contributions of the $2{:}1$ splits cancel.
Nevertheless, exact enumeration shows that the three estimators have proportional mean drifts at $K=3$: the mean baseline gives exactly $\tfrac23F$ and GRPO exactly $\sqrt2F$ (\cref{app:drift}). \emph{For $K=3$, the choice of estimator changes the time scale but not the shape of $F$}.
For larger $K$ the shapes differ. Each $K$ and estimator then has its own exact drift $F_K$, window edge $z^\star_K$ (where $F_K'=0$) and edge drift $F^\star_K=F_K(z^\star_K)$.
\Cref{tab:window} (Appendix) lists $z^\star_K$, which does not increase with $K$, so larger groups narrow the unstable window.
The analysis below uses $F$ only through the sign of $F'$ and the ratio $F^\star/F$, so it holds for any $K$ and estimator once $F$, $z^\star$ and $F^\star$ are replaced by $F_K$, $z^\star_K$ and $F^\star_K$. We quote $K=3$ values and drop the subscript unless stated otherwise.



% ---------------------------------------------------------------------------
\section{Collective Theory}
\label{sec:theory}

We now couple many problems through shared parameters. We introduce a minimal model, show that a balanced population amplifies shared perturbations while its problems are undecided, derive the critical number of problems from a race between two time scales, and extend the result to heterogeneous confidence and many answers. We then give a kernel form of the theory for neural networks and show how to stabilize the uniform collective mode.

\subsection{A minimal many-problem model}
Consider $N$ binary problems whose logits share one parameter:
\begin{equation}
  z_i=u_i+\lambda v,\qquad i=1,\dots,N.
  \label{eq:model}
\end{equation}
$u_i$ is problem-specific, $v$ is a shared direction, and $\lambda$ measures how strongly each problem's logit loads on it.
In a language model, $v$ stands for anything that moves all answers the same way: a format preference, a bias toward short answers, a preferred numerical value, or the output bias of an answer token.
Training uses minibatch SGD with step $\eta$. Each step draws $B$ problems and averages their contributions, $\theta\leftarrow\theta+\frac{\eta}{B}\sum_{i\in\text{batch}}\hat g_i\,\nabla_\theta z_i$, where $\hat g_i$ is the estimate of \cref{sec:single} for problem $i$ and $\theta=(u_1,\dots,u_N,v)$. Since $\partial z_i/\partial u_j=\delta_{ij}$ and $\partial z_i/\partial v=\lambda$, the step moves $u_i$ by $\frac{\eta}{B}\hat g_i$ and $v$ by $\frac{\eta\lambda}{B}\sum_{i\in\text{batch}}\hat g_i$.
A problem is in the batch with probability $B/N$, and $\mathbb E[\hat g_i]=F(z_i)$, so the expected step does not depend on $B$. The batch size affects only the noise.
Replacing the stochastic update by this expected drift, which is accurate for small steps, gives
\begin{align}
  \dot u_i&=\tfrac{\eta}{N}F(z_i), &
  \dot v&=\tfrac{\eta\lambda}{N}\textstyle\sum_jF(z_j),\nonumber\\
  \Rightarrow\ \dot z_i&=\underbrace{\tfrac{\eta}{N}F(z_i)}_{\text{local sharpening}}+\underbrace{\eta\lambda^2\langle F\rangle}_{\text{shared drift}},
  \label{eq:mf}
\end{align}
where $\langle\cdot\rangle$ averages over problems.
We call \cref{eq:mf} the mean-field equations. Because the problems interact only through $v$, which sums all their gradients, the average $\langle F\rangle$ enters exactly. The only approximation is dropping the sampling noise. Without noise a perfectly balanced population would stay balanced forever, so noise re-enters below only as the seed $\delta$ of the shared mode.
The two terms scale differently with $N$. Each problem receives a $1/N$ share of the training signal, but the shared parameter aggregates the signal of all problems.

\subsection{Shared perturbations grow while problems are undecided}
We first ask whether a balanced population, in which different problems prefer different answers, resists a common push.
Suppose the initial preferences are balanced: half the problems lean to $+$ and half to $-$, with $|z_i(0)|=a$.
The constant $a>0$ is the \emph{initial confidence}, the logit magnitude shared by all problems. By symmetry the problems stay at $\pm a(t)$ as they sharpen, with $a(0)=a$.
By oddness of $F$, $\langle F\rangle=0$, so $v\equiv0$ is a solution and each problem sharpens independently.
Perturb it by a common logit shift $w=\lambda\,\delta v$. Linearizing \cref{eq:mf} gives the following (derivation in \cref{app:stability}).

\begin{proposition}[Shared-mode instability]
\label{prop:instab}
To leading order in $1/N$, a common shift evolves as $\dot w=\eta\lambda^2\langle F'(z(t))\rangle\,w$, with the logits $z(t)$ evaluated along the sharpening trajectory.
A small shared perturbation therefore grows at time $t$ if and only if $\langle F'(z(t))\rangle>0$. For a homogeneous population this means $a(t)<z^\star$.
\end{proposition}

\Cref{prop:instab} is a statement about the \emph{instantaneous} growth rate of a perturbation along a time-dependent trajectory, on which every $|z_i|$ keeps increasing. It is not the instability condition of a fixed point, and it is not sufficient for collapse: a shared perturbation that grows for a while can still be too small to flip any answer by the time sharpening has carried the population out of the window. Whether collapse occurs is decided by the \emph{integrated} growth, which we compute in \cref{sec:race}.

The structure is that of the paramagnetic instability of the Curie--Weiss model \citep{nishimori2001}.
$w$ plays the role of a magnetization, the problems are spins coupled only through the mean field, and $\langle F'\rangle$ is the inverse of a susceptibility that changes sign.
Undecided problems ($|z|<z^\star$) collectively amplify \emph{any} shared bias, whether it comes from sampling noise or a slight initial asymmetry.
Confident problems ($|z|>z^\star$) collectively suppress it.

\subsection{A race between sharpening and the shared mode}
\label{sec:race}
\Cref{prop:instab} says when the shared mode grows, but not whether it grows enough to change any answer.
The instability is transient, because local sharpening moves every problem out of the window.
The outcome is decided by which process wins.
Each problem sharpens as in \cref{eq:sharp}, slowed by the factor $1/N$: $\dot a=\frac{\eta}{N}F(a)$. The total log-amplification of the shared mode while the population is inside the window is
\begin{equation}
  \Lambda=\int\!\eta\lambda^2F'(a(t))\,dt=\lambda^2N\,G(a),
  \label{eq:Lambda}
\end{equation}
with $G(a)=\int_a^{z^\star}(F'/F)\,dx=\ln[F^\star/F(a)]$ (\cref{app:stability}).
Local sharpening takes time $\propto N$, and the shared mode grows at an $N$-independent rate $\propto\lambda^2$, so a larger training set gives the shared mode more time to grow.
If the seed of the shared mode has size $\delta$, the minority problems are flipped once $\delta e^{\Lambda}=O(1)$, which gives a critical number of problems
\begin{equation}
  \boxed{N_c\simeq\frac{\ln(1/\delta)}{\lambda^2\,\ln[F^\star/F(a)]}}
  \label{eq:Nc}
\end{equation}
Above $N_c$ every problem converges to the same answer. Accuracy drops to the base rate while the majority-vote reward goes to~1.
For other group sizes and estimators, $G$ is replaced by $G_K(a)=\ln[F^\star_K/F_K(a)]$, built from the exact drift of \cref{sec:single}. \Cref{tab:window} lists $G_K(0.5)$, which also does not increase with $K$. A constant factor in $F$ cancels in the ratio, so $G_3=G$ for all three estimators.

\Cref{eq:Nc} leaves the seed $\delta$ unspecified.
In stochastic training $\delta$ comes from minibatch sampling noise injected into $v$.
Let $\sigma$ be the standard deviation of the single-problem estimate $\hat g_i$, which fluctuates because each update draws fresh answers. Averaging over a batch of $B$ problems gives per-step noise $\eta\lambda^2\sigma/\sqrt B$ on the common shift. With this noise accumulated over the growth time $1/(\eta\lambda^2F')$, we get $\delta^2\sim\eta\lambda^2\sigma^2/(BF')$ (\cref{app:seed}).
The seed enters $N_c$ only through $\ln(1/\delta)=\tfrac12\ln B+\ln(1/\lambda)+\text{const}$, so batch size has only a logarithmic effect. Under the linear scaling rule $\eta\propto B$ it has no effect at all.

\paragraph{Falsifiable predictions.}
Together, these results give four predictions, which we test in \cref{sec:sims}.
(P1)~A sharp transition in $N$ emerges from the dynamics and is not imposed.
(P2)~$N_c\propto\lambda^{-2}$, up to the logarithmic correction above.
(P3)~Batch size enters only logarithmically.
(P4)~In the homogeneous fixed-coupling model, if $a>z^\star$, a small shared perturbation decays at every instant, for every $N$ and every training length. This is a stability property, not a time-scale effect. It is a prediction for this model only: in a network whose kernel changes during training, a small initial $G$ does not guarantee stability (\cref{sec:nn}).
A corollary is that higher-entropy (less confident) initial models have more problems inside the window and therefore collapse at smaller $N$, which matches the empirical observation of \citet{he2026urlvr}.

\subsection{Generalizations}
\label{sec:general}
The analysis so far assumes a homogeneous population of binary problems. We now relax both assumptions.

For arbitrary initial logits, \cref{eq:Lambda} becomes $\Lambda=\lambda^2N\,G[\rho_0]$ (\cref{app:stability}), where
\begin{equation*}
  G[\rho_0]=\sup_s\int_0^s\langle F'(z(s'))\rangle\,ds'
\end{equation*}
and each $z$ follows the sharpening equation \cref{eq:sharp}, $dz/ds=F(z)$ in rescaled time $s=\eta t/N$, from its initial value.
The integrand is the instability index $\langle F'\rangle$ of \cref{prop:instab}, evaluated along the sharpening trajectory.

Next we drop the restriction to two answers.
With $M$ candidate answers, each problem has softmax logits $z_i=u_i+\lambda v\in\mathbb R^M$, and the shared parameter $v\in\mathbb R^M$ is a bias toward some answers over others.
Let $g(z)\in\mathbb R^M$ be the expected single-problem update, the analogue of $F$.
In the symmetric state each problem prefers one answer by a margin $a$, and the preferred answers are spread evenly over the $M$ labels.
A common shift $w=\lambda\,\delta v$ then evolves as $\dot w=\eta\lambda^2\bar Jw$, where $\bar J$ is the Jacobian of $g$ averaged over problems.
For a single problem the Jacobian singles out its preferred answer and couples the components of $w$.
Averaging over problems removes this coupling: every label is preferred by the same number of problems, so $\bar J$ treats all labels alike.
Adding the same constant to all logits leaves a softmax unchanged, so only the $M-1$ zero-sum directions of $w$ matter, and on them $\bar J=\kappa(a)\,I_{M-1}$ (\cref{app:multi}).
The shared mode therefore splits into $M-1$ decoupled directions, each obeying $\dot w=\eta\lambda^2\kappa(a)\,w$.
Each is the binary problem of \cref{prop:instab} with $F'$ replaced by $\kappa$.
Likewise the margin $a$ between the preferred and the other logits sharpens as $\dot a=\frac\eta Nf(a)$. Here $f=g_s-g_{\neq s}$ is the drift of the logit of the preferred answer $s$ minus that of any other answer, and it plays the role of $F$.
The race of \cref{sec:theory} carries over unchanged. The window is $a<a^\star$ with $\kappa(a^\star)=0$, and $\Lambda=\lambda^2N\,G_M(a)$ with $G_M(a)=\int_a^{a^\star}\kappa(x)/f(x)\,dx$.
Both $\kappa$ and $f$ follow from exact enumeration over multinomial count vectors (\cref{app:drift,app:multi}). For $M=2$ they reduce to $F'$ and $F$.

\subsection{Neural networks: a kernel form}
\label{sec:kernel}
The results so far keep the coupling of \cref{eq:model}, a single shared direction with a fixed strength $\lambda$. A real network instead couples problems through all of its parameters.
For a general network, linearizing the parameter dynamics gives $\dot z=\frac{\eta}{N}\mathbf K F(z)$, where $\mathbf K_{ij}=\nabla_\theta z_i\!\cdot\!\nabla_\theta z_j$ is the empirical neural tangent kernel \citep{jacot2018ntk}. $\mathbf K_{ij}$ measures how much a gradient step on problem $j$ moves the logit of problem $i$.
Two collective directions matter. One is the uniform mode $\mathbf 1$, which is the shared direction above. The other is the \emph{sign mode} $s$, with $s_i=\mathrm{sign}\,z_i(0)$ the answer problem $i$ currently prefers.
Moving along $\mathbf 1$ shifts every problem toward the same answer, which is collapse. Moving along $s$ makes every problem more confident in its own answer, which is sharpening applied to all problems at once.
A kernel component $k_s s_is_j$ means that a step sharpening problem $j$ also pushes $z_i$ toward $s_i$, so sharpening one problem sharpens the others. This is how a network generalizes, and it competes with the uniform mode.
We summarize $\mathbf K$ by four numbers.
The \emph{self-kernel} $k_d$, the mean of $\mathbf K_{ii}$, is how much a step on a problem moves its own logit, that is, the effective step size.
The \emph{shared coupling} $k_u$, the mean of $\mathbf K_{ij}$ over $i\ne j$, is how much a step on one problem moves every other logit the same way, whatever their answers. It drives collapse and plays the role of $\lambda^2$.
The \emph{prediction-aligned coupling} $k_s$ is the mean of $s_is_j\mathbf K_{ij}$ over $i\ne j$. It protects against collapse.
The remainder $k_d'=k_d-k_u-k_s$ is the purely problem-specific part of the self-kernel and sets the rate of local sharpening.
Projecting $\mathbf K\approx k_d'\mathbf I+k_u\mathbf 1\mathbf 1^\top+k_s ss^\top$ onto these modes and repeating the race of \cref{sec:race} yields (\cref{app:kernel})
\begin{equation}
  \Lambda(N)=G\,\frac{N\,k_u}{k_d'+N\,k_s},\qquad
  N_c=\frac{\Lambda_c\,k_d'}{G\,k_u-\Lambda_c\,k_s}.
  \label{eq:kernel}
\end{equation}
The projection keeps exactly the parts of $\mathbf K$ that enter the race: its diagonal and its action along $\mathbf 1$ and $s$, since $\mathbf 1^\top\mathbf K\mathbf 1/N\approx k_d'+Nk_u$ and $s^\top\mathbf K s/N\approx k_d'+Nk_s$.
It is exact in the toy model, which is the special case $k_d'=1$, $k_u=\lambda^2$, $k_s=0$. In a network it drops the other collective modes of $\mathbf K$, which we return to in \cref{app:nnremedy}.
When $k_s>0$ the growth of $\Lambda$ with $N$ saturates at $Gk_u/k_s$, and there is a \emph{feature-protected} regime $G\,k_u<\Lambda_c\,k_s$ in which no training-set size causes collapse.
All quantities in \cref{eq:kernel} can be measured on the initial model before any RL step.
This assumes that $\mathbf K$ stays near its initial value, as in the lazy regime. When the network learns features, its gradients $\nabla_\theta z_i$ and hence $\mathbf K$ change, and the prediction can fail. Re-measuring $\mathbf K$ later does not help once the race is decided, and after collapse $k_s=k_u$ because all $s_i$ agree.
In the toy model, a feature coupling $z_i\mathrel{+}=\mu s_iw$ realizes $k_s=\mu^2$ and gives $N_c=1/(\lambda^2G/\Lambda_c-\mu^2)$.

\subsection{Stabilizing the uniform collective mode}
\label{sec:remedy}
Finally, we turn from diagnosing the instability to damping it.
In the minimal model the instability lives entirely in the uniform shared mode, so it can be targeted without slowing per-problem sharpening. A network has further collective modes, and damping the uniform one alone does not protect accuracy there (\cref{sec:nnresults}).
We add to each minibatch objective a marginal-entropy bonus $\beta\,H(\bar p)$ with weight $\beta\ge0$. Here $\bar p=\frac1B\sum_{i\in\text{batch}}p_i$ is the probability of answer $+$ averaged over the batch, and $H$ is the binary entropy
\begin{equation*}
  H(q)=-q\ln q-(1-q)\ln(1-q).
\end{equation*}
The majority-vote reward lowers the entropy of each problem's answer, and the bonus raises the entropy of the answers pooled across problems, in the spirit of mutual-information objectives for unsupervised classification \citep{bridle1992phantom,krause2010rim}.
When answers are balanced, $\bar p=\tfrac12$ is the maximum of $H$, so the bonus exerts no force and per-problem sharpening is unaffected.
A common shift moves every $p_i$ the same way and pushes $\bar p$ away from $\tfrac12$, and the bonus pushes back.
Near the symmetric state it contributes $-4\beta\langle p(1-p)\rangle^2$ to the growth rate of the shared mode (\cref{app:remedy}), so the shared mode is damped when
\begin{equation}
  \beta>\beta_c(a)=\frac{F'(a)}{4\,[p(1-p)]^2}\le\tfrac12 ,
  \label{eq:beta}
\end{equation}
with $p=\sigma(a)$. The threshold is largest at $a=0$, where $F'(0)=1/8$ and $p(1-p)=\tfrac14$ give exactly $\tfrac12$, so any $\beta>\tfrac12$ damps the uniform shared mode at every confidence.
